\documentclass[../DoAn.tex]{subfiles}
\begin{document}

\begin{center}
    \Large{\textbf{TÓM TẮT NỘI DUNG ĐỒ ÁN}}\\
\end{center}
\vspace{1cm}

Trong bối cảnh thiết bị bay không người lái (Drone) ngày càng được sử dụng rộng rãi trong nhiều lĩnh vực, nhu cầu phát hiện và cảnh báo sớm drone đóng vai trò quan trọng trong bảo đảm an ninh, an toàn hạ tầng và quản lý không phận tầm thấp. So với các phương pháp dựa trên camera, âm thanh hoặc radar, phát hiện thụ động bằng tín hiệu vô tuyến (Radio Frequency - RF) có ưu điểm không phụ thuộc điều kiện quan sát, ít chịu ảnh hưởng bởi điều kiện môi trường và có khả năng khai thác trực tiếp tín hiệu liên lạc giữa drone và bộ điều khiển.

Đồ án nghiên cứu bài toán phát hiện drone từ tín hiệu RF dưới dạng phân loại nhị phân giữa hai lớp \textit{Drone} và \textit{Non-drone}. Tín hiệu RF được biểu diễn dưới dạng tín hiệu I/Q (In-phase/Quadrature), sau đó được chuyển đổi thành biểu diễn thời gian–tần số dưới dạng ảnh phổ (spectrogram) bằng phép biến đổi Fourier thời gian ngắn (Short-Time Fourier Transform - STFT), sau đó được chuẩn hóa và sử dụng làm đầu vào cho các mô hình học sâu. Hai chiến lược khai thác bộ trích xuất đặc trưng (backbone) tiền huấn luyện được khảo sát, gồm phân loại dựa trên trích xuất đặc trưng (\textit{Feature Extraction-based Classification}) và tinh chỉnh mô hình (\textit{fine-tuning}), nhằm đánh giá hiệu quả của từng phương pháp trên nhiều kiến trúc CNN (Convolutional Neural Network – mạng nơ-ron tích chập) và kiến trúc Transformer.

Các thí nghiệm được thực hiện trên tập dữ liệu huấn luyện xây dựng từ dữ liệu RF công khai kết hợp với dữ liệu RF tự thu, đồng thời đánh giá bổ sung trên các tập dữ liệu ngoài miền nhằm kiểm chứng khả năng tổng quát hóa của mô hình trong các điều kiện thu khác nhau. Hiệu năng của các mô hình được phân tích thông qua các chỉ số đánh giá phân loại và khả năng duy trì kết quả trên những tập dữ liệu chưa được sử dụng trong quá trình huấn luyện. Kết quả thực nghiệm cho thấy \textit{fine-tuning} cải thiện rõ rệt hiệu năng trên dữ liệu trong miền, trong khi đánh giá ngoài miền cho thấy hiệu quả triển khai còn phụ thuộc vào backbone, chiến lược huấn luyện và điều kiện thu tín hiệu. Những kết quả này cho thấy tính khả thi của phương pháp phát hiện drone dựa trên tín hiệu RF, đồng thời nhấn mạnh tầm quan trọng của đánh giá ngoài miền đối với các hệ thống hướng tới triển khai thực tế.


\end{document}
